\section*{第 \S\arabic{exercisegroup} 节习题}
\phantomsection
\addcontentsline{toc}{section}{第 \protect\S\arabic{exercisegroup} 节习题}

\begin{enumerate}
\def\labelenumi{\arabic{enumi})}
\tightlist
\item
  设 \(\mathbf{f}\) 是一个 \(n\) 次可导的向量函数（在区间 \(\mathrm{I} \subset \mathbf{R}\) 上）。证明公式
\end{enumerate}

\[
\mathrm{D} ^ {n} \mathbf {f} \left(e ^ {x}\right) = \sum_ {m = 1} ^ {n} \frac {a _ {m}}{m !} e ^ {m x} \mathbf {f} ^ {(m)} \left(e ^ {x}\right)
\]

在使 \(e^{x} \in I\) 的每一点 x 处成立，其中系数 \(a_{m}\) 可表为

\[
a _ {m} = \sum_ {p = 0} ^ {m} (- 1) ^ {p} \binom {m} {p} (m - p) ^ {n}
\]

（用 I, p.~41, exerc. 7 的方法，并利用 \(e^{x}\) 的泰勒展开）。

\begin{enumerate}
\def\labelenumi{\arabic{enumi})}
\setcounter{enumi}{1}
\item
  设 \(f\) 是一个实函数，在一点 \(n\) 次可导，该点为 \(x\) ；\(\mathbf{g}\) 是一个向量函数，\(n\) 次可导，可导点为 \(f(x)\) 。若令 \(\mathrm{D}^n(\mathbf{g}(f(x))) = \sum_{k=1}^{n} \mathbf{g}^{(k)}(f(x)) u_k(x)\) ，则 \(u_k\) 仅依赖于函数 \(f\) ；由此推出：\(u_k(x)\) 是 \(t^k\) 的系数（在 \(e^{-tf(x)}\mathrm{D}^n(e^{tf(x)})\) 按 \(t\) 展开的展开式中）。
\item
  对每个实数 \(x > 0\) 与每个复数 \(m = \mu + iv\) ，令 \(x^{m} = e^{m\log x}\) ；试证：III, p.~108 的公式 (19) 对复的 \(m\) 与 \(x > -1\) 仍然有效，且该公式中的余项 \(r_{n}(x)\) 满足不等式
\end{enumerate}

\[
\begin{array}{l l} | r _ {n} (x) | \leqslant \left| \binom {m - 1} {n} \frac {m}{\mu} x ^ {n} \Big [ (1 + x) ^ {\mu} - 1 \Big ] \right| & \text { if } \mu \neq 0 \\ | r _ {n} (x) | \leqslant \left| \binom {m - 1} {n} m x ^ {n} \log (1 + x) \right| & \text { if } \mu = 0. \end{array}
\]

把二项级数收敛性的研究推广到 m 为复数的情形。

\begin{enumerate}
\def\labelenumi{\arabic{enumi})}
\setcounter{enumi}{3}
\tightlist
\item
  对每个实数 \(x\) 及每个数 \(p > 1\) ，证明不等式
\end{enumerate}

\[
\begin{array}{r l} | 1 + x | ^ {p} \leqslant 1 + p x + \frac {p (p - 1)}{2} x ^ {2} + \dots + \frac {p (p - 1) \dots (p - m + 2)}{(m - 1) !} x ^ {m - 1} \\ & + \frac {p (p - 1) \dots (p - m + 1)}{m !} | x | ^ {m} + h _ {p} | x | ^ {p} \end{array}
\]

其中令 \(m = [p]\) （即 \(p\) 的整数部分），并且

\[
h _ {p} = \frac {p (p - 1) \dots (p - m + 1)}{(m - 1) !} \int_ {0} ^ {1} z ^ {p - m} (1 - z) ^ {m - 1} d z.
\]

¶ 5) 试证：\(\pi\) 是无理数，方法如下：若 \(\pi = p/q\) （ \(p\) 与 \(q\) 为整数），则令 \(f(x) = (x(\pi - x))^n / n!\) 后，积分 \(q^n \int_0^\pi f(x) \sin x dx\) 将是一个整数 \(>0\) （利用 \(n + 1\) 阶分部积分公式）；但试证另一方面，\(q^n \int_0^\pi f(x) \sin x dx\) 当 \(n\) 趋于 \(+\infty\) 时趋于 0 。

\begin{enumerate}
\def\labelenumi{\arabic{enumi})}
\setcounter{enumi}{5}
\tightlist
\item
  试证：在区间 \([-1,+1]\) 上，函数 \(|x|\) 是多项式的一致极限，方法是注意 \(|x|=(1-(1-x^{2}))^{1/2}\) 并利用二项级数。由此得出魏尔斯特拉斯定理的另一证明（cf.~II, p.~83, exerc. 20）。
\end{enumerate}

¶ 7) 设 \(p\) 是素数，\(\mathbf{Q}_p\) 是 \(p\) 进数域（Gen.~Top., III, p.~322 and 323, exercs. 23 to 25），\(\mathbf{Z}_p\) 是 \(p\) 进整数环，\(\mathfrak{p}\) 是主理想 \((p)\) （在 \(\mathbf{Z}_p\) 中）。

\begin{enumerate}
\def\labelenumi{\alph{enumi})}
\tightlist
\item
  设 \(a = 1 + pb\) ，其中 \(b \in \mathbf{Z}_p\) 是乘法群 \(1 + \mathfrak{p}\) 的元素；试证：当有理整数 \(m\) 无限增大时， \(p\) 进数 \(\frac{(1 + pb)^{p^m} - 1}{p^m}\) 趋于一个极限，它等于收敛级数
\end{enumerate}

\[
\frac {p b}{1} - \frac {p ^ {2} b ^ {2}}{2} + \dots + (- 1) ^ {n - 1} \frac {p ^ {n} b ^ {n}}{n} + \dots
\]

把这个极限记作 \(\log a\) 。

\begin{enumerate}
\def\labelenumi{\alph{enumi})}
\setcounter{enumi}{1}
\item
  试证：当 \(p\) 进数 \(x\) 在 \(\mathbf{Q}_p\) 中趋于 0 时，数 \(\frac{a^x - 1}{x}\) 趋于 \(\log a\) （利用 \(a\) ) 及 \(\mathbf{Q}_p\) 的拓扑的定义）。
\item
  试证：若 \(p \neq 2\) ，则 \(\log a \equiv pb \pmod{\mathfrak{p}^2}\) ；若 \(p = 2\) ，则 \(\log a \equiv 0 \pmod{\mathfrak{p}^2}\) ，且 \(\log a \equiv -4b^4 \pmod{\mathfrak{p}^3}\) 。
\item
  试证：若 \(p \neq 2\) （相应地，p = 2），则 \(x \mapsto \log x\) 是乘法拓扑群 \(1 + p\) 到加法拓扑群 p 上的同构（相应地，到 \(p^{2}\) 上）；特别地，若 \(e_{p}\) 是 \(1 + p\) 中使 \(\log e_{p} = p\) 的元素（相应地，\(\log e_{2} = 4\) ），则 \(Z_{p}\) 到 \(1 + p\) 上的、与 \(x \mapsto \frac{1}{p} \log x\) 互逆的同构（相应地，与 \(x \mapsto \frac{1}{4} \log x\) 互逆者）就是 \(y \mapsto e_{p}^{y}\) （cf.~Gen.~Top., III, p.~323, exerc. 25）。
\item
  试证：对每个 \(a \in 1 + p\) ，连续函数 \(x \mapsto a^{x}\) （定义在 \(Z_{p}\) 上）在每一点处的导数等于 \(a^{x} \log a\) ；由此推出函数 \(\log x\) 在 \(1 + p\) 的每一点处的导数等于 1/x 。
\end{enumerate}

¶ 8) a) 用习题 7 的记号，试证：通项为 \(x^n / n!\) 的级数对每个 \(x \in \mathfrak{p}\) 收敛（若 \(p \neq 2\) ），对每个 \(x \in \mathfrak{p}^2\) （但对任何 \(x \notin \mathfrak{p}^2\) 都不）收敛（若 \(p = 2\) ）（定出 \(p\) 在 \(n!\) 的素因子分解中的指数）。若 \(f(x)\) 是这级数的和，试证 \(f\) 是 \(\mathfrak{p}\) （相应地，\(\mathfrak{p}^2\) ）到 \(1 + \mathfrak{p}^2\) 中的连续同态。由此推出：对每个 \(z \in \mathbf{Z}_p\) ，有 \(f(pz) = e_p^z\) （相应地，\(f(p^2z) = e_p^z\) ）；换言之，

\[
e _ {p} = 1 + \frac {p}{1 !} + \frac {p ^ {2}}{2 !} + \dots + \frac {p ^ {n}}{n !} + \dots \quad \text { if } p \neq 2
\]

\[
e _ {2} = 1 + \frac {4}{1 !} + \frac {4 ^ {2}}{2 !} + \dots + \frac {4 ^ {n}}{n !} + \dots
\]

（利用习题 7 e)）。

\begin{enumerate}
\def\labelenumi{\alph{enumi})}
\setcounter{enumi}{1}
\tightlist
\item
  对每个 \(a \in 1 + \mathfrak{p}\) 及每个 \(x \in \mathbf{Z}_p\) ，试证 \(\log (a^x) = x\log a\) ，并由 a) 与习题 7 d) 推出
\end{enumerate}

\[
a ^ {x} = 1 + \frac {x \log a}{1 !} + \frac {x ^ {2} (\log a) ^ {2}}{2 !} + \dots + \frac {x ^ {n} (\log a) ^ {n}}{n !} + \dots
\]

\begin{enumerate}
\def\labelenumi{\alph{enumi})}
\setcounter{enumi}{2}
\tightlist
\item
  对每个 \(m \in Z_{p}\) ，试证：连续函数 \(x \mapsto x^{m}\) （定义在 \(1 + p\) 上）的导数等于 \(mx^{m-1}\) （利用 b) 与习题 7 e)）。
\end{enumerate}

\(d\) ) 试证：对每个 \(m\in \mathbf{Z}_p\) 及每个 \(x\in \mathfrak{p}\) （当 \(p\neq 2\) ）或每个 \(x\in \mathfrak{p}^2\) （当 \(p = 2\) ），通项为 \(\binom{m}{n}x^n\) 的级数收敛，且其和是 \(m\) 的连续函数；由此推出这个和等于 \((1 + x)^m\) ，方法是注意 \(\mathbf{Z}\) 在 \(\mathbf{Z}_p\) 中（处处）稠密。

¶ 9) 用习题 7) 的记号，把 \(\mathbf{O}_2^+ (\mathbf{Q}_p)\) （空间 \(\mathbf{Q}_p^2\) 的旋转群）定义为形如

\[
\left( \begin{array}{c c} x & y \\ - y & x \end{array} \right)
\]

的矩阵所成的群，其元素属于 \(\mathbf{Q}_p\) ，满足 \(x^2 + y^2 = 1\) ，该群赋有 Gen.~Top., VIII, p.~125, exerc. 2 中定义的拓扑。

\begin{enumerate}
\def\labelenumi{\alph{enumi})}
\item
  用 \(G_{n}\) 表示 \(\mathbf{O}_{2}^{+}(\mathbf{Q}_{p})\) 中由满足 \(t=y/(1+x)\in\mathfrak{p}^{n}\) 的矩阵所成的子群。试证：\(G_{n}\) 是紧群，\(G_{n}/G_{n+1}\) 同构于 Z/pZ，且 \(G_{1}\) 的仅有的紧子群就是诸群 \(G_{n}\) （cf.~Gen.~Top., III, p.~323, exerc. 24）。
\item
  试证：\(G_{1}\) 与由矩阵
\end{enumerate}

\[
\left( \begin{array}{c c} x & y \\ - y & x \end{array} \right)
\]

所成的子群相同，其中 \(x^{2} + y^{2} = 1\) ，且 \(x \in 1 + p^{2}\) ， \(y \in p\) （若 \(p \neq 2\) ）；\(x \in 1 + p^{3}\) ， \(y \in p^{2}\) （若 p = 2）。

\begin{enumerate}
\def\labelenumi{\alph{enumi})}
\setcounter{enumi}{2}
\tightlist
\item
  试证：通项为 \((-1)^{n} x^{2n}/(2n)!\) 与 \((-1)^{n-1} x^{2n+1}/(2n+1)!\) 的级数对每个 \(x \in p\) 收敛（若 \(\neq 2\) ），对每个 \(x \in p^{2}\) 收敛（若 p = 2）；设 \(\cos x\) 与 \(\sin x\) 是这两个级数的和。试证映射
\end{enumerate}

\[
x \mapsto \left( \begin{array}{c c} \cos x & \sin x \\ - \sin x & \cos x \end{array} \right)
\]

是加法拓扑群 \(\mathfrak{p}\) （相应地，\(\mathfrak{p}^2\) ）到群 \(G_1\) 上的同构。

\(d\) ) 若 \(p\) 形如 \(4h + 1\) （ \(h\) 为整数），则 \(\mathbf{Q}_p\) 中存在元素 \(i\) 使得 \(i^2 = -1\) 。对每个 \(z \in \mathbf{Q}_p^*\) 关联矩阵

\[
\left( \begin{array}{c c} \frac {1}{2} \left(z + \frac {1}{z}\right) & \frac {1}{2 i} \left(z - \frac {1}{z}\right) \\ - \frac {1}{2 i} \left(z - \frac {1}{z}\right) & \frac {1}{2} \left(z + \frac {1}{z}\right) \end{array} \right)
\]

便定义了乘法群 \(Q_{p}^{*}\) 到群 \(\mathbf{O}_{2}^{+}(\mathbf{Q}_{p})\) 上的一个同构；在这个同构下，群 \(1 + p\) 对应于 \(G_{1}\) ，并且有 \(\cos px + i \sin px = e_{p}^{ix}\) （III, p.~126, exerc. 8）。

\begin{enumerate}
\def\labelenumi{\alph{enumi})}
\setcounter{enumi}{4}
\tightlist
\item
  若 p 形如 \(4h + 3\) （h 为整数），则 \(\mathbf{O}_{2}^{+}(\mathbf{Q}_{p})\) 的矩阵的元素必为 p 进整数。此时多项式 \(X^{2} + 1\) 在 \(Q_{p}\) 中不可约；设 \(\mathbf{Q}_{p}(i)\) 是 \(Q_{p}\) 的一个二次扩张，由添加 \(X^{2} + 1\) 的一个根 i 得到；给 \(\mathbf{Q}_{p}(i)\) 赋以 Gen.~Top., VIII, p.~127, exerc. 2 中定义的拓扑。群 \(\mathbf{O}_{2}^{+}(\mathbf{Q}_{p})\) 同构于 \(\mathbf{Q}_{p}(i)\) 中范数为 1 的元素所成的乘法群 N，同构的方式是把矩阵 \(\begin{pmatrix} x & y \\ -y & x \end{pmatrix}\) 对应于元素 \(z = x + iy\) 。试证：有 \(p + 1\) 个数满足方程 \(x^{p+1} = 1\) （在 \(\mathbf{Q}_{p}(i)\) 中），且它们构成 N 的一个循环子群 R（论证方法同 Gen.~Top., III, p.~323, exerc. 24；然后证明有 \(p + 1\) 个不同的根满足同余式 \(x^{p+1} \equiv 1 \pmod{p}\) （在 \(\mathbf{Q}_{p}(i)\) 中），并对该同余式的每个根 a 作序列 \((a^{p2n})\) ）。由此推出：群 \(\mathbf{O}_{2}^{+}(\mathbf{Q}_{p})\) 同构于群 R 与 \(G_{1}\) 的乘积。
\end{enumerate}

\(f\) ) 试证：当 \(p = 2\) 时，群 \(\mathbf{O}_2^+ (\mathbf{Q}_p)\) 同构于群 \(\mathrm{G}_1\) 与 4 阶循环群的乘积。
% semantic-fragments: legacy-input-seam
\relax{}
